\section{Toward a Structural ToE Master Equation}
\label{sec:toe-master}
%% ============================================================

We now compress the previously established structural requirements into a
single master equation. This equation is not introduced as an external ansatz,
but as the minimal common form forced by the triolic structure, the
Heisenberg compensator, the probe--closure system, the projected closure
residue, and the phase-shifted observer frame.

\begin{equation}\label{eq:toe_master_structural}
\mathfrak G[\Psi,g]
+\mathfrak T_{\mathrm{tri}}[\Psi]
+\mathfrak H_{\mathrm{HC}}[\Psi]
+\mathfrak Q_{\mathrm{cl}}[\Psi]
+\mathfrak B_{\beta}[\Psi]
=0 .
\end{equation}

Here $\mathfrak G[\Psi,g]$ denotes the geometric sector,
$\mathfrak T_{\mathrm{tri}}[\Psi]$ the triolic sector dynamics,
$\mathfrak H_{\mathrm{HC}}[\Psi]$ the Heisenberg-compensator contribution,
$\mathfrak Q_{\mathrm{cl}}[\Psi]$ the closure term determined by the projected
closure residue, and $\mathfrak B_{\beta}[\Psi]$ the phase-shifted observer
term. Equation \eqref{eq:toe_master_structural} is therefore not an
independent postulate, but the structural compression of the results obtained
in the preceding sections.

For a semi-concrete realization, one may write
\begin{equation}\label{eq:toe_master_semiconcrete}
\mathbf{R}[\Psi]
+\sum_{a=1}^{3} D_a\Psi_a
+\sum_{a\neq b}\Lambda_{ab}\Psi_b
+\lambda_{\mathrm{HC}}(I-C)\Psi
+\lambda_{\mathrm{cl}}\mathcal R_{\mathrm{cl}}[\Psi]
+\lambda_{\beta}(\Psi-U_{\beta}\Psi)
=0 .
\end{equation}
Here $\mathbf{R}[\Psi]$ denotes the effective geometric contribution,
$D_a$ are the operators of the triolic decomposition,
$\Lambda_{ab}$ encode inter-sector couplings,
$C$ is the Heisenberg compensator,
$\mathcal R_{\mathrm{cl}}[\Psi]$ is the projected closure residue,
and $U_{\beta}$ is the phase-shift operator of the observer frame.
Equation \eqref{eq:toe_master_semiconcrete} should be read only as a
semi-concrete representative of the structural master equation
\eqref{eq:toe_master_structural}.

\begin{theorem}[Structural necessity of the ToE master equation]
\label{thm:toe_master_necessity}
Assume that the theory is required to satisfy simultaneously:
\begin{enumerate}
\item a common underlying field structure $\Psi$;
\item a triolic decomposition
\[
\Psi=\Psi_1\oplus\Psi_2\oplus\Psi_3;
\]
\item a dual unitary symmetry between the first two sectors;
\item a Heisenberg compensator selecting the collapse/light output;
\item a probe--closure system with projected closure residue;
\item a phase-shifted observer frame;
\item an Einstein limit on large-scale closure-stable configurations.
\end{enumerate}
Then the dynamics cannot be described by a purely geometric, purely sectorial,
or purely observer-based equation alone. Instead, it must take the form of a
single structural field equation
\[
\mathfrak G[\Psi,g]
+\mathfrak T_{\mathrm{tri}}[\Psi]
+\mathfrak H_{\mathrm{HC}}[\Psi]
+\mathfrak Q_{\mathrm{cl}}[\Psi]
+\mathfrak B_{\beta}[\Psi]
=0,
\]
containing at least these five structural blocks.
\end{theorem}

\begin{proof}
Each structural requirement excludes an incomplete description.

A purely geometric equation does not encode the triolic decomposition, the
Heisenberg compensator, or the phase-shifted observer frame. A purely
triolic-sector equation does not account for the emergent geometric reading of
the closure-stable background. A purely observer-based equation does not
select physically admissible configurations. The projected closure residue
shows that closure must enter dynamically, and the probe--closure system shows
that this closure is not external to the field content itself. Finally, the
Einstein limit demonstrates that the geometric sector is not optional, but a
distinguished large-scale projection of the same structure.

Hence none of the five structural contributions can be removed without
violating one of the established premises. Therefore the dynamics must be
governed by a single combined equation containing at least these five
structural blocks.
\end{proof}

\begin{theorem}[Uniqueness up to structural equivalence]
\label{thm:structural-uniqueness}
Let $\Psi$ and $\widetilde\Psi$ be two nontrivial solutions realizing the
structural core and satisfying the master equation
\eqref{eq:toe_master_structural}. Then they belong to the same structural
class up to sector-preserving unitary equivalence, compensator projection, and
observer-frame phase shift.
\end{theorem}

\begin{proof}
Both solutions are constrained by the same triolic decomposition, the same
compensator logic, the same projected closure requirement, and the same
observer-frame mechanism. Their admissible differences can therefore only lie
in the choice of representative within one common structural organization:
sector exchange is absorbed by unitary symmetry, collapse/light output is fixed
up to compensator projection, and effective readout is changed only by the
phase-shifted observer frame. Once these structurally allowed transformations
are factored out, no independent second theory remains. Thus the admissible
solutions are unique up to structural equivalence.
\end{proof}

\begin{proposition}[From local triolic structure to global field dynamics]
\label{prop:local-to-global}
The transition from local triolic structure to global field dynamics is forced
by the following chain:
\[
\begin{aligned}
&\text{local triolic decomposition}
  \longrightarrow \text{duality and compensator action}
  \longrightarrow \text{projected closure}\\
&\text{closure-stable background geometry}
  \longrightarrow \text{global master equation}.
\end{aligned}
\]
In particular, Dirac-, QFT-, and Einstein-type regimes arise as structural
projections of one common global dynamics rather than as separate starting
points.
\end{proposition}

\begin{proof}
Locally, the theory resolves into three sectors linked by duality and
compensator projection. The projected closure residue then decides whether
these local sectors assemble into a physically admissible configuration.
Whenever closure is achieved, the resulting sector can be read as a stable
background supporting geometric reconstruction, effective matter dynamics, and
observer-frame projection. The global master equation is therefore not added
from outside, but is the minimal common dynamics of all local closure-stable
triolic realizations. Its distinguished restrictions yield the Dirac-, QFT-,
and Einstein-type regimes established in the preceding sections.
\end{proof}

\begin{corollary}[Triolic realization of admissible solutions]
\label{cor:toe_triolic_realization}
Let $\Psi$ satisfy the structural master equation
\eqref{eq:toe_master_structural}. Then every physically admissible solution
decomposes into three necessary sectors,
\[
\Psi=\Psi_1\oplus\Psi_2\oplus\Psi_3,
\]
where $\Psi_1$ is the matter sector, $\Psi_2$ its dual tachyonic sector, and
$\Psi_3$ the collapse/light sector selected by the Heisenberg compensator.
\end{corollary}

\begin{proof}
The triolic term enforces the three-sector structure, the compensator term
links the first two sectors and selects the collapse/light output, and the
closure term excludes non-admissible configurations by the projected closure
residue. Thus the three sectors are not optional additions, but the necessary
realization of admissible solutions.
\end{proof}

\begin{corollary}[Structural projections of the master equation]
\label{cor:toe_structural_projections}
The master equation \eqref{eq:toe_master_structural} admits the following
distinguished structural projections:
\begin{enumerate}
\item in the probe regime, it induces the effective Dirac-type equation of the
probe--closure system;
\item in the phase-shifted observer regime, it yields an effective QFT-type
description after compensator decomposition;
\item in the large-scale closure-stable regime, it reduces to the sharpened
Einstein limit.
\end{enumerate}
\end{corollary}

\begin{proof}
Each projection is obtained by restricting one and the same structural
equation to a distinguished regime.

In the probe regime, the auxiliary probe sector couples to the background
through the same closure and compensator structure, yielding the effective
Dirac-type operator. In the observer regime, the phase-shifted frame and the
compensator decomposition produce the effective field description seen by a
matter observer. In the large-scale regime, the closure residue and the
sectorial fluctuations are suppressed, so that the geometric contribution
dominates and yields the sharpened Einstein limit.

Hence probe dynamics, effective QFT, and gravitational geometry are not
separate theories, but structural projections of one common closure-stable
equation.
\end{proof}

The significance of the master equation is that it does not add a new layer to
the theory, but reveals the common body already implicit in the triolic
structure, the Heisenberg compensator, the probe--closure system, and the
closure analysis. Matter, tachyon, photon, probe dynamics, effective QFT, and
the Einstein limit now appear as sectorial realizations or limiting
projections of one and the same closure-stable structural equation.

\section{Structural Reading of Matter, Gauge, Mass, and Higgs}
\label{sec:structural-reading-physics}
%% ============================================================

The structural picture may now be read in compact physical form. Matter is not
an externally added ingredient, but the closure-stable bound sector of the
triolic decomposition. The tachyonic sector is its dual partner under the
unitary exchange symmetry, while the photon/light sector is the
Heisenberg-compensator output selected by collapse and closure. In this sense,
matter, tachyon, and photon are not separate ontological primitives, but
sectorial realizations of one common structural body.

Gauge structure is read off from the admissible transitions that preserve
closure and compensator consistency. What appears as gauge symmetry in the
effective observer description is therefore the residual invariance of
closure-compatible sector transport. In the same way, mass is not introduced
as an independent primitive parameter, but as a closure/frame residue: the
effective inertia seen by a matter observer records the nontrivial cost of
maintaining a closure-stable bound sector under phase-shifted readout.

This motivates the following compressed structural reading.

\begin{proposition}[Structural reading of matter, gauge, and mass]
\label{prop:structural-matter-gauge-mass}
Within the closure-stable triolic system determined by
\eqref{eq:toe_master_structural}:
\begin{enumerate}
\item matter is the closure-stable bound realization of the first sector;
\item the photon/light sector is the compensator-selected collapse output;
\item effective gauge symmetry is the invariance of closure-compatible
sector transport;
\item effective mass is the observer-frame residue of a nontrivial
closure-stable bound configuration.
\end{enumerate}
\end{proposition}

\begin{proof}
The first two statements summarize the triolic decomposition together with the
centrality of the Heisenberg compensator. Gauge structure was already shown to
arise from the admissible symmetry content of the effective sector picture;
within the present structural closure framework this means precisely that the
allowed sector transitions must preserve closure and compensator consistency.
Finally, the phase-shifted observer frame and the probe--closure analysis show
that nontrivial bound sectors are not read out as pure closure identities but
as effective residues. These residues are perceived as inertial or mass-like
parameters by the observer frame.
\end{proof}

The Higgs question should then not be read as the introduction of an isolated
extra object, but as the identification of the structural locus at which
coupling, mediation, and effective mass readout are organized. At the present
level of the theory, the safest statement is therefore not a final derivation,
but a structural hypothesis.

\begin{remark}[Structural role of the Higgs sector]
\label{rem:structural-higgs-role}
The Higgs sector is expected to arise either as the effective coupling locus of
the closure-stable triolic system or as the observer-frame image of a deeper
mediating point within the septinion/meta-septum organization. In either
reading, Higgs physics is not external to the structural core, but a special
readout of the same closure, compensator, and frame-shift mechanism.
\end{remark}

%% ============================================================



\section{Structural Normalization for the Next Companion Stage}
\label{sec:structural-normalization}

\begin{remark}[Version note]
Version \docversion\ introduces a structural normalization layer for the
Companion. The Millennium matrix is read primarily as the global matrix
realization of the Heisenberg compensator; the tool apparatus is split into
general Millennium-matrix tools and special Heisenberg-compensator tools; the
carrier/group/physical-readout distinction is made explicit; and the global
reader orientation is organized through a master chain from triolic source to
field-equation readout.
\end{remark}

\begin{remark}[Null axis versus light axis]
In structurally prior parts of the Companion we use the term \emph{null axis}
for the distinguished compensated axis selected by the MM/HC architecture.
When the same axis is read at the physical propagation level, it may be called
the \emph{light axis}. Thus ``light axis'' is not a second object, but the
physical readout of the structurally prior null axis.
\end{remark}

\section{The Millennium Matrix as Global Heisenberg-Compensator}
\label{sec:millennium-hc}

\begin{definition}[Millennium matrix]
The Millennium matrix $\mathcal M_{16}$ is the global structural matrix that
transports one and the same Sphaera-object through its admissible algebraic,
closure, and projection layers. It preserves identity through viewpoint
change.
\end{definition}

\begin{definition}[Layer decomposition]
We write
\[
\mathcal M_{16}
=
\mathcal M_{\mathrm{tri}}
+
\mathcal M_{\mathrm{alg}}
+
\mathcal M_{\mathrm{cl}}
+
\mathcal M_{\mathrm{proj}},
\]
where $\mathcal M_{\mathrm{tri}}$ encodes the triolic kernel organization,
$\mathcal M_{\mathrm{alg}}$ admissible algebraic realization changes,
$\mathcal M_{\mathrm{cl}}$ closure compatibility, and
$\mathcal M_{\mathrm{proj}}$ admissible final readout.
\end{definition}

\begin{definition}[Global Heisenberg-compensator reading]
The Millennium matrix is said to be read as a global Heisenberg-compensator
matrix if its action extends the local rule
\[
C=\tfrac12(\mathrm{id}+J)
\]
to the full multi-layer architecture of the Companion. Under this reading, the
Millennium matrix is denoted by
\[
\mathcal H_{16},
\]
or identified with $\mathcal M_{16}$ when no ambiguity arises.
\end{definition}

\begin{theorem}[Structural role of the Millennium matrix]
\label{thm:mm-role}
The Millennium matrix is the global structural intermediary of the Companion in
the following sense:
\begin{enumerate}
  \item it preserves the identity of the triolic core under admissible
  realization change;
  \item it transports structural content through admissible algebraic layers
  without ontology change;
  \item it enforces closure discipline before readout;
  \item when read as $\mathcal H_{16}$, it becomes the global balancing
  mechanism for physically admissible interpretation.
\end{enumerate}
\end{theorem}

\begin{proof}
The triolic layer supplies the primitive object. The admissible algebraic
layers supply realizations of that object. Closure selects the admissible
sector, while projection only reads out what closure has stabilized. The
Millennium matrix is the unique global carrier of this ordered transport.
Under the Heisenberg-compensator reading, the same transport becomes the
global balancing mechanism across all admissible layers.
\end{proof}

\section{Carrier, Group, and Physical Readout Layers}
\label{sec:carrier-group-readout}

\begin{definition}[Carrier layer]
The carrier layer is the total admissible structural support on which the
Sphaera architecture is realized. It contains the triolic source structure,
the admissible algebraic realizations, the closure discipline, and the
transport induced by the Millennium matrix.
\end{definition}

\begin{definition}[Group layer]
The group layer is the symmetry layer of the Companion. It consists of the
admissible group actions, generators, and representation changes acting on
already given carrier-layer structure.
\end{definition}

\begin{definition}[Physical readout layer]
The physical readout layer is the final effective layer in which already
stabilized structural and symmetry data are interpreted as physically
meaningful sectors, observables, and field equations.
\end{definition}

\begin{theorem}[Order of interpretation]
\label{thm:carrier-group-readout-order}
The admissible order of interpretation is
\[
\text{carrier layer}
\;\longrightarrow\;
\text{group layer}
\;\longrightarrow\;
\text{physical readout layer}.
\]
\end{theorem}

\begin{proof}
Without a carrier there is nothing for a symmetry group to act on. Symmetries
therefore organize admissible realizations of carrier-layer content. Only
after structural transport and symmetry organization have been fixed may a
physical interpretation be attached.
\end{proof}

\section{Master Structural Theorem of the Companion}
\label{sec:master-structural-theorem}


\begin{theorem}[Master structural theorem]
\label{thm:master-structural}
The Companion is governed by one common structural architecture with the
following ordered logic:
\[
\begin{aligned}
&\text{triolic source}
 \;\longrightarrow\;
 \text{admissible carrier}
 \;\longrightarrow\;
 \mathcal M_{16}
 \;\longrightarrow\;
 \mathcal H_{16}
 \;\longrightarrow\;
 \text{closure-stable sector}\\[0.2em]
&\longrightarrow\;\text{group-layer symmetry readout}
 \;\longrightarrow\;\text{null axis}\\[0.2em]
&\longrightarrow\;\text{photon fix-sector}
 \;\longrightarrow\;\text{residual localization}\\[0.2em]
&\longrightarrow\;\text{matter / tachyon / photon sector readout}\\[0.2em]
&\longrightarrow\;\text{Higgs mediation}
 \;\longrightarrow\;\text{field-equation readout}.
\end{aligned}
\]
\end{theorem}

\begin{proof}
The triolic source supplies the primitive structural object. The carrier layer
provides its admissible support, while the Millennium matrix transports this
same object through admissible realization changes. Under the
Heisenberg-compensator reading, the same matrix becomes the global balancing
mechanism that stabilizes the closure-compatible sector. The group layer then
supplies admissible symmetry actions on already given stabilized structure.
The null axis is the structurally distinguished sector carried through the
cascade, and its terminal compensated readout is the photon fix-sector.
Relative to that fix-sector, closure-stable residual localization yields the
matter branch, while its $J$-dual realization yields the tachyon branch. The
Higgs sector mediates between the compensated photon reference and the
persistence of residual localization. Finally, the effective field equations
read out the dynamical content of this already stabilized architecture.
\end{proof}

\begin{remark}[Compact roadmap for the reader]
The reader may keep the following compact roadmap in mind:
\[
\boxed{\begin{aligned}
&\text{one source} \to \text{one transported architecture}\\
&\to \text{one compensated closure logic}\\
&\to \text{three principal sector readouts} \to \text{one mediated field dynamics}
\end{aligned}}
\]
\end{remark}

\subsection{Schematic Master Diagram}
\label{subsec:master-diagram}

\begin{definition}[Schematic master diagram]
\label{def:master-diagram}
The structural logic of the Companion may be compressed into the following
schematic master diagram:
\[
\mathfrak T
\;\longrightarrow\;
\mathfrak C
\;\xrightarrow{\;\mathcal M_{16}\;}
\mathfrak C^{\mathrm{adm}}
\;\xrightarrow{\;\mathcal H_{16}\;}
\mathfrak C^{\mathrm{cl}}
\;\longrightarrow\;
\mathfrak G
\;\longrightarrow\;
\mathfrak L
\;\longrightarrow\;
\mathfrak P
\;\longrightarrow\;
\mathfrak R
\;\longrightarrow\;
\mathfrak H
\;\longrightarrow\;
\mathfrak F.
\]
Here:
\begin{align*}
\mathfrak T &:= \text{triolic source},\\
\mathfrak C &:= \text{carrier layer},\\
\mathfrak C^{\mathrm{adm}} &:= \text{admissibly transported carrier},\\
\mathfrak C^{\mathrm{cl}} &:= \text{closure-stable compensated carrier},\\
\mathfrak G &:= \text{group-layer symmetry readout},\\
\mathfrak L &:= \text{null-axis structure},\\
\mathfrak P &:= \text{photon fix-sector},\\
\mathfrak R &:= \text{residual localization sector},\\
\mathfrak H &:= \text{Higgs-mediation sector},\\
\mathfrak F &:= \text{field-equation readout}.
\end{align*}
\end{definition}

\begin{remark}[Language regime along the master chain]
\label{rem:language-regime-master-chain}
The Companion follows a structure-first logic along the master chain up to the
point at which a stable physical readout becomes available. Accordingly, the
earlier stages
\[
\mathfrak T
\to
\mathfrak C
\to
\mathcal M_{16}
\to
\mathcal H_{16}
\to
\mathfrak G
\to
\mathfrak L
\]
are described primarily in structurally prior language. By contrast, once the
readout has stabilized at the level of
\[
\mathfrak P
\to
\mathfrak R
\to
\mathfrak H
\to
\mathfrak F,
\]
physically explicit language is not only allowed but required. Thus the shift
from structural terminology to physical terminology is intentional and marks a
readout threshold, not a change of underlying object.
\end{remark}

\subsection{Matter, Tachyon, and Photon as Three Readouts of One Branch Architecture}
\label{subsec:mtp-readout}

\begin{definition}[Matter branch]
The matter branch is the closure-stable bound branch whose physical readout
retains nontrivial localization relative to the photon fix-sector.
\end{definition}

\begin{definition}[Tachyon branch]
The tachyon branch is the $J$-dual readout of the matter branch:
\[
K_2 = J\!\cdot\!K_1.
\]
It is the branch in which the same structural source is read in its frame-dual,
non-bound, or overextended transport mode relative to the compensated photon
sector.
\end{definition}

\begin{definition}[Photon branch]
The photon branch is the canonical compensated fix-sector readout $K_3$, i.e. the terminal null-axis sector selected by the Heisenberg compensator.
\end{definition}

\begin{theorem}[Three-readout principle]
\label{thm:three-readout-principle}
Matter, tachyon, and photon are three structurally linked readouts of one
common branch architecture: matter is the bound closure-stable localization
readout, tachyon is the $J$-dual unbound readout of the same branch source,
and photon is the compensated fix-sector readout.
\end{theorem}

\begin{proof}
Admissible sector transport does not create new ontology. The compensated
sector is the canonical fix-sector. Hence the matter branch is determined by
persistent closure-stable localization relative to this sector. The involution
$J$ generates the frame-dual branch of the same source, which yields the
tachyonic readout. The compensator selects the stabilized null/fix-sector,
which yields the photon readout. Therefore all three branches belong to one
common structural architecture and differ only by their mode of realization and
readout.
\end{proof}

\subsection{Mass Readout Relative to the Null Axis and Photon Fix-Sector}
\label{subsec:mass-null-axis}

\begin{definition}[Photon fix-sector revisited]
The photon sector is the canonical compensated readout of the global null
axis:
\[
L=\bigcap_{n\ge 0} L_n = S_0.
\]
It is the sector in which the frame-sensitive imbalance has been fully
neutralized and the structural readout lies entirely in the HC-fix sector.
\end{definition}

\begin{definition}[Massive branch]
A branch is called massive if it remains closure-stably localized and does not
collapse completely into the global null/fix sector. Equivalently, a massive
branch is one whose physical readout retains a nontrivial localization residue
relative to the photon sector.
\end{definition}

\begin{theorem}[Massive versus massless readout]
\label{thm:massive-vs-massless}
Within the MM/HC architecture, the distinction between massive and massless
physical readout is the distinction between complete reduction to the
compensated photon fix-sector and closure-stable persistence of localization
relative to that fix-sector.
\end{theorem}

\begin{proof}
If a branch reduces completely to the null/fix sector, no residual
localization remains, and the physical readout is massless. If, however, a
branch remains closure-stably localized and does not reduce completely to the
null/fix sector, then a nonzero localization residue survives relative to the
photon reference sector. This is precisely the massive readout.
\end{proof}

\subsection{Higgs Readout as Mediation Between Fix-Sector and Residual Localization}
\label{subsec:higgs-fix-localization}

\begin{definition}[Higgs mediation principle]
\label{def:higgs-mediation}
The Higgs readout is the effective mediation principle that governs how
closure-stable localization is maintained, transferred, or expressed relative
to the compensated photon fix-sector.
\end{definition}

\begin{theorem}[Structural role of the Higgs sector]
\label{thm:higgs-role-refined}
Within the MM/HC architecture, the Higgs sector is the distinguished coupling
readout that mediates between the compensated photon/null fix-sector and the
effective mass readout of closure-stable localized branches.
\end{theorem}

\begin{proof}
By the previous null-axis analysis, the photon sector is the terminal
compensated fix-sector. By the mass analysis, massive readout is residual
localization relative to that fix-sector. Hence the remaining structural task
is to identify the admissible interface through which this residual
localization is expressed in the physical readout layer. The Higgs sector is
precisely this interface.
\end{proof}

\section{Field Equations as Physical Readout of the MM/HC Structure}
\label{sec:field-readout}

\begin{definition}[Field-readout principle]
\label{def:field-readout-principle}
A field equation is admissible in the present framework if it is obtained as
the effective physical readout of carrier-layer structure, group-layer symmetry
organization, and MM/HC-mediated closure stabilization.
\end{definition}

\begin{definition}[Minimal structural field content]
The minimal physical readout of the closure-stable MM/HC architecture is
carried by the tuple
\[
(\Psi,\mathcal A_\mu,\hat P_{C_7},\mathcal H_C),
\]
where $\Psi$ is the effective matter/probe field, $\mathcal A_\mu$ the
effective symmetry connection induced by the group layer,
$\hat P_{C_7}$ the closure projector selecting the admissible stable sector,
and $\mathcal H_C$ the effective Heisenberg-compensator field.
\end{definition}

\begin{theorem}[Minimal MM/HC field equation]
\label{thm:minimal-mmhc-field-equation}
The minimal local field equation compatible with the carrier layer, the group
layer, and the MM/HC closure architecture is
\[
\bigl(i\Gamma^\mu\mathcal D_\mu-\Omega\bigr)\Psi=0,
\]
with effective operator
\[
\Omega
=
m_0\mathbf 1
+\alpha\,\hat P_{C_7}\mathcal M_{16}\hat P_{C_7}
+\beta\,\mathcal T
+\gamma\,\mathcal H_C.
\]
\end{theorem}

\begin{proof}
The carrier layer fixes the underlying admissible object and its triolic
organization. The group layer supplies the symmetry transport term encoded by
$\mathcal A_\mu$. The MM/HC architecture contributes closure-stable internal
transport via $\hat P_{C_7}\mathcal M_{16}\hat P_{C_7}$ and compensator
balancing via $\mathcal H_C$. Finally, the triolic source remains visible at
the level of effective coupling, which yields the term $\mathcal T$. No term
may be removed without losing one previously established layer.
\end{proof}
